\chapter[The Physical Massless Multiplier]{The Massless Fermionic Interaction on the Physical Hilbert Space}
\label{ch:massless-multiplier}

\section{From a finite Fock operator to multiplication on wavefunctions}
Let \(M\) be the original paired \(\mathrm{SU}(2)\) algebra data,
\(X\in\mathcal X\) the bosonic configuration, and \(\mathcal F\)
the original Fock space. The source massless frequency
\(K_{1,0}(X)\) equals \(K_{1,1}(X)-K_{1,1}(0)\) in the
source conventions. This difference removes the constant mass term.
Write \(B(X)=M.\operatorname{quantizeFrequency}(K_{1,0}(X))\).
Its restriction of scalars to the reals is exactly
\(M.\operatorname{deformedFermionField}(1,0,X)\). In particular
the real inner product in the source energy identity is the real
part of the complex Fock pairing with \(B(X)\). This uses the
source normalization of the Majoranas and the sign established
by the massless quantization proof.
\formalrecord{M1a supplies the complex quantizer, M1b the field
identity, and M1c the coefficient difference. The
\href{https://github.com/etblink/math/actions/runs/45059474635}{target integration receipt}
includes their exact source closure.}

Each \(B(X)\) acts on the finite Fock space. Its dependence on \(X\)
leads to a separate operator \(V\) on
\(\mathcal H=L^2(\mathcal X;\mathcal F)\):
\[
 D(V)=\{u\in\mathcal H:B(\cdot)u(\cdot)\in L^2\},
 \qquad (Vu)(X)=B(X)u(X)\quad\text{a.e.}
\]
Here the membership condition is on a representative and is independent
of that choice. Two representatives agree outside a null set, where
their images under the same pointwise linear map also agree.
Continuity of the source field in \(X\) supplies measurability and
local boundedness. Smooth compactly supported vectors belong to
\(D(V)\). Their density in the ambient space gives density of \(D(V)\).

To see the graph argument, suppose \(u_n\to u\) and \(Vu_n\to v\) in
\(L^2\). Extract a joint subsequence converging almost everywhere in
both coordinates. On that common full-measure set, continuity of the
finite linear map \(B(X)\) in its input gives
\(v(X)=B(X)u(X)\). Since \(v\in L^2\), \(u\in D(V)\) and \(Vu=v\).
The source proof separates the maximal domain, the algebraic graph,
the almost-everywhere limit argument, and the paired specialization.
\formalrecord{D0a through D0e1d supply these interfaces. The generic
partial operator was independently qualified in
\href{https://github.com/etblink/math/actions/runs/45039242640}{D0c0}.
The paired closed dense graph is covered by the
\href{https://github.com/etblink/math/actions/runs/45059474635}{unified target receipt}.}

\section{Restriction to the original physical fixed space}
For the original gauge action \(U_g\), source covariance preserves
\(D(V)\) and gives \(VU_gu=U_gVu\) there. Thus a fixed input has a
fixed output. Let \(E=\mathcal H_{\rm phys}\) and write
\(V_{\rm phys}\) for the restriction with domain \(D(V)\cap E\)
and values in \(E\). If \(i:E\hookrightarrow\mathcal H\) is the
inclusion, its actual graph satisfies
\[
 \operatorname{Graph}(V_{\rm phys})
   =(i\times i)^{-1}\bigl(\operatorname{Graph}(V)\bigr).
\]
Continuity of \(i\times i\) transfers closedness. Separately, the
original invariant smooth core is dense in \(E\) and belongs to the
physical multiplier domain. These arguments give the closed, densely
defined source operator. A self-adjoint realization would require
an adjoint-domain argument for this operator.
\formalrecord{GM0--GM3 supply gauge covariance, domain invariance and
intertwining. GM4A constructs the restriction. GM4B1 supplies density
and \href{https://github.com/etblink/math/actions/runs/45085548204}{GM4B2 R2}
checks the inverse-image graph and paired closedness.}

\section{The first composition with the sixteen-charge form}
\label{sec:multiplier-energy-join}
The October FCP S4B34--36 chain already established the concrete
signed Pauli/Fock identity and evaluated the historical closed
charge energy on its embedded invariant smooth core. Its hypotheses
and input domain belong to that tested construction. The present
unified source has the same mathematical ingredients for a focused
composition with the named physical multiplier.

For an original invariant smooth vector \(f\), let \(u=i_{\rm core}f\)
be its image in \(E\). The source graph embedding and the massless
core-domain theorem give \(u\in D(T)\cap D(V_{\rm phys})\). The
following is the proposed direct consumer of the existing
interfaces:
\begin{align}
 \tfrac1{16}\|Tu\|^2
 &=\tfrac12\sum_p\|\iota\partial_p f\|_{L^2}^2
   +\int_{\mathcal X} V_B(X)\|f(X)\|^2\,dX
   +\operatorname{Re}\langle u,V_{\rm phys}u\rangle_E .
 \label{eq:proposed-multiplier-energy}
\end{align}
The ordinary composition is short. Restrict the charge-form
identity to \(u=i_{\rm core}f\), then replace the original pointwise
massless fermion field by the real restriction of \(B(X)\). The
almost-everywhere action of \(V\) and the isometric inclusion of
\(E\) identify its integrated pairing with the final term. The
Cauchy--Schwarz inequality gives integrability of this pairing from
\(u,Vu\in L^2\). This manuscript records the composition and its
source dependencies. An exact, separately accepted Lean consumer
of the displayed statement remains a research target.

A later extension should name the intended form domain and prove a
continuity or approximation estimate there. A graph approximation
for \(T\) alone controls the charge norm. Convergence in that norm
does not provide convergence in the graph norm of \(V_{\rm phys}\).
The question of which energy terms extend individually therefore
requires analytic input about the proposed domain. The
source comparison and bounded formal target are listed in
\texttt{EDITION\_0\_9\_SOURCE\_COMPARISON.md}.
